% TOP_PROBLEM: {"id": 30006956, "title": "Fujita Very Ampleness Conjecture", "category_id": 5, "category": {"id": 5, "name": "algebraic_geometry", "display_name": "Algebraic Geometry", "description": "Geometric objects defined by polynomial equations.", "slug": "algebraic-geometry", "order_index": 5, "created_at": "2026-07-31T15:26:25.671Z"}, "status": "open"}
% FIRST_POSTED: 2026-09-27
% !TeX program = lualatex
% ATTEMPT_STATUS: unresolved
% Model: GPT-6 Astra Ultra; continuation dated 27 September 2026.
\documentclass[11pt]{article}
\usepackage[margin=25mm]{geometry}
\usepackage{fontspec}
\setmainfont{Latin Modern Roman}
\usepackage{amsmath,amssymb,mathtools}
\usepackage{unicode-math}
\setmathfont{Latin Modern Math}
\usepackage{xurl,hyperref}
\hypersetup{colorlinks=true,urlcolor=blue}
\setcounter{secnumdepth}{0}
\setlength{\parindent}{0pt}
\setlength{\parskip}{5pt}
\setlength{\emergencystretch}{3em}
\begin{document}
\section*{\texorpdfstring{Fujita Very Ampleness Conjecture}{Fujita Very Ampleness Conjecture}}\label{30006956}
\subsection{Definitions and mathematical statement}
Let $X$ be a smooth complex projective variety of dimension $n$, and let $L$ be an ample Cartier divisor. A line bundle is very ample when its complete linear system defines a closed immersion into projective space. With $K_X$ the canonical divisor, Fujita's assertion here is
\[
 \mathcal O_X(K_X+(n+2)L)\quad\text{is very ample}.
\]
In particular this requires enough global sections not only to avoid base points but also to separate distinct points and tangent directions. The integer $n+2$ is fixed uniformly by dimension. The freeness conjecture with $n+1$ is a different, weaker property and is listed separately in Section~\href{https://mehmetmars7.github.io/Erdosproblems-llm-hunter/problem.html?type=open_problems&id=problem.fujita-freeness-conjecture}{363}. Replacing $L$ by a rational divisor without an integrality convention would not define the same line bundle.
\par\smallskip\textit{Formulation and concept references:} \href{https://arxiv.org/abs/2408.06530}{[S1]}.

\subsection{Short English statement}
Does adding n plus two copies of any ample divisor to the canonical divisor always yield a projective embedding of an n-dimensional smooth variety?
\subsection{Research attempt}

\paragraph{The sharp target and the source audit (27 September 2026).}
Write line bundles additively and put $M=K_X+(n+2)L$. The target is
very ampleness for this exact coefficient and every ample integral $L$.
It is not enough to construct one nonzero section, remove base points,
obtain a birational map, or replace $n+2$ by a larger unspecified integer.
Murayama's primary text [S1] distinguishes all these issues and records
the very-ampleness conjecture in characteristic zero as known through
dimension two. It also records a stronger established subcase than the
one proved below, namely when $L$ is ample and globally generated.

The current-source search located Siu's 2026 article [E1], whose abstract
claims uniform very-ampleness and jet-generation bounds with larger
coefficients. That is not the sharp $n+2$ assertion. This attempt does
not treat the newer bound as a solution of the printed conjecture,
and does not reproduce or independently audit that analytic proof.

Our positive results are complete proofs for curves, for very ample
$L$ in every dimension using Kodaira vanishing, and for explicit
products of curves. A blow-up calculation then isolates a sufficient
local-positivity condition and shows where the general ample case remains.
The standard inputs used below are Serre duality on curves, Kodaira
vanishing, adjunction, Bertini's theorem, and the basic cohomology of
blowing up smooth points. Their use is explicit; no effective global
positivity theorem at the conjectural bound is assumed.

\paragraph{Length-two subschemes locate the actual obstruction.}
For a zero-dimensional subscheme $Z\subset X$ of length two, tensor
its ideal sequence by $M$:
\[
 0\longrightarrow M\otimes\mathcal I_Z\longrightarrow M
 \longrightarrow M|_Z\longrightarrow0.
\]
Here $h^0(Z,M|_Z)=2$. Surjectivity of
\[
 H^0(X,M)\longrightarrow H^0(Z,M|_Z)
 \tag{466.1}
\]
for all such $Z$ is the usual very-ampleness criterion on a smooth
projective variety. A reduced $Z=\{p,q\}$ tests separation of points.
A nonreduced length-two scheme supported at $p$ tests one tangent
direction at $p$. These tests also imply global generation by choosing
length-two schemes containing each point. The resulting proper morphism
separates points and has injective differential; equivalently its
complete linear system is a closed immersion.

Since $(n+2)L$ is ample, Kodaira vanishing gives $H^1(X,M)=0$ when
$n\geq1$. The long exact sequence therefore gives the particularly
precise equivalence
\[
 \text{(466.1) is onto}\quad\Longleftrightarrow\quad
 H^1(X,M\otimes\mathcal I_Z)=0.
 \tag{466.2}
\]
Kodaira vanishing for $M$ alone does not prove the vanishing with
$\mathcal I_Z$. The ideal sheaf is not an ample line bundle, and
this extra condition is exactly the interpolation problem.

\paragraph{The complete curve case.}
\textbf{Proposition 466.1.} If $C$ is a smooth projective complex curve
and $L$ is ample, then $K_C+3L$ is very ample.

\emph{Proof.}
Let $d=\deg L\geq1$ and let $Z$ have length two. On a smooth curve
$Z$ is an effective Cartier divisor of degree two, including the case
$Z=2p$. Serre duality gives
\[
 H^1(C,K_C+3L-Z)^*\cong H^0(C,Z-3L).
\]
The degree on the right is $2-3d<0$, so there is no nonzero section:
a nonzero section would define an effective divisor of that negative
degree. Hence the left side vanishes and (466.1) is surjective for
every $Z$.\hfill$\square$

The same argument proves that $K_C+2L$ is globally generated. For a
single point $p$, the dual obstruction has degree $1-2d<0$.
This latter statement will be the base of an induction, not a substitute
for the preceding length-two calculation.

The coefficient three is necessary uniformly. For $C=\mathbb P^1$
and $L=\mathcal O(1)$, one has $K_C+2L=\mathcal O$, which defines
a constant map and cannot embed the curve, whereas
$K_C+3L=\mathcal O(1)$. Large genus does not spoil the proof:
$\deg(K_C+3L)=2g-2+3d\geq2g+1$, and the argument uses the precise
degree of the dual obstruction rather than a heuristic abundance of sections.

\paragraph{A tensor-product lemma for preserving an embedding.}
\textbf{Lemma 466.2.} If $A$ is very ample and $B$ is globally generated
on a smooth complex projective variety, then $A+B$ is very ample.

\emph{Proof.}
Fix a length-two $Z$. Choose a section $b\in H^0(X,B)$ nonzero at
every point of its support. For one support point this follows from
global generation. For two points, the sections vanishing at either
point form two proper linear subspaces of $H^0(X,B)$; over the infinite
field $\mathbb C$ their union is not the whole space.
The restriction $b|_Z$ is a unit in a trivialization of $B|_Z$, including
when $Z$ is nonreduced, since it is nonzero in the residue field.
Multiplication by $b|_Z$ is thus an isomorphism from $H^0(Z,A|_Z)$
to $H^0(Z,(A+B)|_Z)$. Sections of $A$ surject onto the former space,
and their products with $b$ are global sections of $A+B$. This proves
the length-two criterion.\hfill$\square$

This lemma requires the factor $A$ to be very ample. Merely ample
$A$ does not automatically separate the two jets, which is why the
lemma does not turn Fujita freeness into the full conjecture by itself.

\paragraph{A proof when the given polarization is very ample.}
\textbf{Proposition 466.3.} For a smooth complex projective $n$-fold
and a very ample $L$, the line bundle $K_X+(n+1)L$ is globally
generated, and $K_X+(n+2)L$ is very ample.

\emph{Proof.}
We prove the global-generation assertion by induction on $n$.
The curve case was proved above. Fix $n\geq2$ and $p\in X$.
Bertini's theorem for hyperplanes through the embedded point $p$
gives a smooth divisor $D\in|L|$ containing $p$: away from $p$ this
linear system has no base points, and at $p$ a general hyperplane has
nonzero differential on $T_pX$ because $L$ embeds $X$.
The restriction $L_D$ is very ample. By adjunction,
\[
 (K_X+(n+1)L)|_D=K_D+nL_D.
\]
The induction hypothesis says that this restricted line bundle is
globally generated. The restriction sequence is
\[
 0\longrightarrow K_X+nL\longrightarrow K_X+(n+1)L
 \longrightarrow K_D+nL_D\longrightarrow0.
\]
Kodaira vanishing gives $H^1(X,K_X+nL)=0$, so its map on global
sections to the right-hand term is onto. A section on $D$ nonzero
at $p$ lifts to a global section on $X$ nonzero there. Since $p$
was arbitrary, global generation follows.

Now apply Lemma 466.2 to $A=L$ and $B=K_X+(n+1)L$.
\hfill$\square$

If desired, a disconnected smooth divisor in this argument can be
handled component by component; global generation on each component
is enough. The substantive condition is smoothness through the chosen
point, which is supplied by the embedding hypothesis. Dimension zero
is immediate and does not affect the induction.

This proof identifies two exact failures for arbitrary ample $L$:
$|L|$ may have no suitable divisor through $p$, and the final tensor
factor $L$ need not separate points or tangent directions. Replacing
$L$ by a very ample multiple $rL$ proves an assertion with coefficient
$(n+2)r$, not with $n+2$. Taking an $r$-th root of a very ample line
bundle is not an operation preserving very ampleness.

\paragraph{Projective-space sharpness and explicit coordinates.}
For $X=\mathbb P^n$, every ample integral divisor is $L=dH$ with
$d\geq1$, and $K_X=-(n+1)H$. Consequently
\[
 K_X+(n+2)L=((n+2)d-n-1)H,
\]
whose coefficient is at least one. To verify the embedding concretely,
the complete linear system of $aH$, $a\geq1$, consists of all degree
$a$ monomials. On the chart $x_0\ne0$, their ratios include
\[
 \frac{x_0^{a-1}x_i}{x_0^a}=\frac{x_i}{x_0}
 \qquad(1\leq i\leq n).
\]
They recover the local coordinates and all tangent directions. The same
holds on every coordinate chart, and pure powers distinguish the
charts; this is the Veronese embedding. In contrast, at $d=1$ the
coefficient $n+1$ yields $K_X+(n+1)L=\mathcal O_X$.
Thus the proposed bound is uniformly sharp in every positive dimension.

The sharpness example is also a warning against numerical shortcuts:
the freeness coefficient may already give global generation with only
a constant map. The step from no base points to an embedding is genuine.

\paragraph{Products of curves with an exterior polarization.}
Let $X=C_1\times\cdots\times C_n$, with each $C_i$ a smooth complex
projective curve, and take
$L=\bigotimes_i p_i^*L_i$, where $\deg L_i=d_i\geq1$.
This is ample. The adjoint line bundle is
\[
 K_X+(n+2)L
   =\bigotimes_i p_i^*(K_{C_i}+(n+2)L_i).
\]
Each factor has degree at least $2g_i+n\geq2g_i+1$, so the same
Serre-duality argument as in Proposition 466.1 makes it very ample.
Taking the product of the resulting embeddings and composing with
the Segre embedding realizes the displayed exterior tensor product as
a very ample line bundle. The product sections give a subsystem
already defining a closed immersion; the full system is therefore
very ample as well.

The qualification on $L$ matters: an arbitrary ample line bundle on a
product need not have this exterior form. This construction establishes
an explicit infinite class of examples without asserting that all
polarizations on all such products have been treated by this argument.

\paragraph{A conditional local-positivity route via blow-ups.}
Assume $n\geq2$. For distinct points $p,q$, let
$\pi:\widetilde X\to X$ be their blow-up with exceptional divisors
$E_p,E_q$. The canonical bundle and ideal transforms satisfy
\[
 K_{\widetilde X}=\pi^*K_X+(n-1)(E_p+E_q),\qquad
 \pi_*\mathcal O_{\widetilde X}(-E_p-E_q)=\mathcal I_{\{p,q\}}.
\]
The higher direct images of this exceptional ideal vanish. More
generally, for the blow-up of a smooth point,
$\pi_*\mathcal O(-rE)=\mathcal I_p^r$ and
$R^j\pi_*\mathcal O(-rE)=0$ for $r\geq0$, $j>0$.
These are standard blow-up identities; on the exceptional projective
space the restrictions are $\mathcal O_{\mathbb P^{n-1}}(r)$,
whose positive-degree cohomology vanishes. They allow the usual Leray
and projection-formula identification of the cohomology groups below.

If the integral divisor
\[
 A_{p,q}=(n+2)\pi^*L-n(E_p+E_q)
 \tag{466.3}
\]
is ample, then
$\pi^*M-E_p-E_q=K_{\widetilde X}+A_{p,q}$.
Kodaira vanishing and those identities yield
$H^1(X,M\otimes\mathcal I_{\{p,q\}})=0$, proving separation of
these two points.

For a single point blow-up, if
\[
 B_p=(n+2)\pi^*L-(n+1)E_p
 \tag{466.4}
\]
is ample, the same argument applied to
$\pi^*M-2E_p=K_{\widetilde X}+B_p$ proves surjectivity onto
$M\otimes\mathcal O_X/\mathfrak m_p^2$. This separates all first
jets simultaneously, and hence every length-two scheme supported at
$p$. Requiring (466.3) for all pairs and (466.4) for all points is
therefore a sufficient condition for the conjectural embedding.

These coefficients are worth recording: two reduced points require
$n$ on each exceptional divisor, whereas the full first-order
neighbourhood requires $n+1$. Confusing the canonical discrepancy
$n-1$ with either of these numbers loses the desired jet condition.

This sufficient condition is far too strong to settle every case.
Already for $X=\mathbb P^n$, $L=H$, and a line through $p,q$, its
strict transform has intersection
\[
 A_{p,q}\cdot\widetilde\ell=(n+2)-2n=2-n.
\]
Thus (466.3) is not ample for $n\geq2$: it has intersection zero
at $n=2$ and negative intersection at $n>2$. Yet the desired adjoint
bundle is $H$ and is very ample. The failure is in this sufficient
blow-up criterion, not in the conjecture. It also shows why obtaining
positivity at two points by one undifferentiated vanishing argument
can waste precisely the amount of positivity that the sharp statement
does not have to spare.

\paragraph{Outcome and first gap.}
The proofs handle the full curve statement, all very ample
polarizations, projective spaces at the sharp bound, and the specified
products of curves. They also produce explicit two-point and jet
vanishing criteria, with a sharp diagnostic showing that those criteria
need not hold even on projective space.

For a general ample $L$, the missing assertion is the vanishing in
(466.2) for every length-two $Z$, or another argument giving the same
evaluation surjectivity without wasting positivity. Neither ordinary
Kodaira vanishing, a very ample multiple, nor an asymptotic result with
a larger coefficient gives that assertion. The sharp conjecture is
unresolved by this attempt.

\subsection{Sources}
\begin{itemize}
\item[S1]Murayama, Moving Seshadri constants and effective Fujita-type conjectures (2024).
\url{https://arxiv.org/abs/2408.06530}.
\textit{Formulation source.}
\item[E1]Yum-Tong Siu, \emph{Very Ampleness Part of Fujita Conjecture
and Multiplier Ideal Sheaves of Higher Order}, Acta Mathematica Sinica,
English Series 42 (2026), 1149--1157, published 15 May 2026.
\url{https://link.springer.com/article/10.1007/s10114-026-6229-8}.
\textit{Publisher abstract and bibliographic record inspected
27 September 2026. Its larger effective bound is distinguished from
the sharp coefficient $n+2$; the proof is not independently audited here.}
\item[E2]Takumi Murayama, full-text HTML of the supplied source,
\url{https://arxiv.org/html/2408.06530}.
\textit{Conjecture 1.1.1, the characteristic-zero and globally-generated
subcases, and the distinction between freeness and jet separation
inspected 27 September 2026.}
\end{itemize}
\end{document}
